% Part: first-order-logic
% Chapter: tableaux
% Section: proving-things

\documentclass[../../../include/open-logic-section]{subfiles}

\begin{document}

\iftag{FOL}
      {\olfileid[fr]{fol}{tab}{pro}}
      {\olfileid[fr]{pl}{tab}{pro}}

\olsection{Exemples de \usetoken{p}{tableau}}

\begin{ex}
Cherchons un !!{tableau} fermé pour l'!!{sentence}
$(!A \land !B) \lif !A$.

Commençons par inscrire l'hypothèse correspondante en haut du
!!{tableau}.
\begin{oltableau}
  [\sFmla{\False}{(\formula{A} \land \formula{B}) \lif \formula{A}},
    just = \TAss]
\end{oltableau}

Il n'y a qu'une hypothèse, donc qu'une seule !!{signed formula} à
laquelle appliquer une règle. Pour toute !!{signed formula}, il existe
toujours au plus une règle applicable : celle qui correspond au signe
et au !!{main operator} de l'!!{sentence}. Dans le cas présent, il faut
donc appliquer $\TRule{\False}{\lif}$.
\begin{oltableau}
  [\sFmla{\False}{(\formula{A} \land \formula{B}) \lif \formula{A}},
    checked, just = \TAss
    [\sFmla{\True}{\formula{A} \land \formula{B}},
      just={\TRule{\False}{\lif}[1]}
      [\sFmla{\False}{\formula{A}}, just={\TRule{\False}{\lif}[1]}]
    ]
  ]
\end{oltableau}
Pour repérer les !!{signed formula}s auxquelles nous avons appliqué la
règle correspondante, nous plaçons une coche à côté de l'énoncé. Il ne
faut cependant placer cette coche que si la règle a été appliquée sur
toutes les branches ouvertes. Dès que la règle correspondante à
!!a{signed formula} a été appliquée sur chacune d'elles, il n'est plus
nécessaire de revenir à cette formule pour appliquer de nouveau la
règle. Il n'y a ici qu'une branche ; une seule application suffit donc.
Les coches ne sont qu'un moyen pratique de construire les tableaux et
ne font pas officiellement partie de leur syntaxe.

Une nouvelle !!{signed formula} permet l'application d'une règle :
$\sFmla{\True}{!A \land !B}$ à la ligne~$2$. L'application de
$\TRule{\True}{\land}$ donne :
\begin{oltableau}
  [\sFmla{\False}{(\formula{A} \land \formula{B}) \lif \formula{A}},
    checked, just = \TAss
    [\sFmla{\True}{\formula{A} \land \formula{B}},
      just={\TRule{\False}{\lif}[1]}, checked
      [\sFmla{\False}{\formula{A}}, just={\TRule{\False}{\lif}[1]}
        [\sFmla{\True}{\formula{A}}, just={\TRule{\True}{\land}[2]}
          [\sFmla{\True}{\formula{B}}, just={\TRule{\True}{\land}[2]}, close
          ]
        ]
      ]
    ]
  ]
\end{oltableau}
La branche contient maintenant $\sFmla{\True}{!A}$ à la ligne~$4$ et
$\sFmla{\False}{!A}$ à la ligne~$3$ ; elle est donc fermée. Comme c'est
la seule branche, le !!{tableau} est fermé. Nous avons ainsi trouvé un
!!{tableau} fermé pour~$(!A \land !B) \lif !A$.
\end{ex}

\begin{ex}
Cherchons maintenant un !!{tableau} fermé pour
$(\lnot !A \lor !B) \lif (!A \lif !B)$.

Commençons par l'hypothèse correspondante :
\begin{oltableau}
  [\sFmla{\False}{(\lnot \formula{A} \lor \formula{B}) \lif
      (\formula{A} \lif \formula{B})}, just=\TAss]
\end{oltableau}
L'unique !!{signed formula} de ce !!{tableau} a pour
!!{main operator}~$\lif$ et pour signe~$\False$. L'application de
$\TRule{\False}{\lif}$ donne :
\begin{oltableau}
  [\sFmla{\False}{(\lnot \formula{A} \lor \formula{B})
      \lif (\formula{A} \lif \formula{B})}, just=\TAss, checked
    [\sFmla{\True}{\lnot \formula{A} \lor \formula{B}},
      just={\TRule{\False}{\lif}[1]}
      [\sFmla{\False}{(\formula{A} \lif \formula{B})},
        just={\TRule{\False}{\lif}[1]}
      ]
    ]
  ]
\end{oltableau}
Nous pouvons maintenant appliquer soit~$\TRule{\True}{\lor}$ à la
ligne~$2$, soit $\TRule{\False}{\lif}$ à la ligne~$3$. L'ordre choisi
est indifférent, pourvu que la règle correspondant à chaque
!!{signed formula} soit appliquée sur toutes les branches. Prenons la
première. La règle $\TRule{\True}{\lor}$ ramifie le !!{tableau} ; ses
deux conclusions deviennent les nouvelles !!{signed formula}s ajoutées
aux deux branches ainsi créées. On obtient :
\begin{oltableau}
  [\sFmla{\False}{(\lnot \formula{A} \lor \formula{B}) \lif
      (\formula{A} \lif \formula{B})}, just=\TAss, checked
    [\sFmla{\True}{\lnot \formula{A} \lor \formula{B}},
      just={\TRule{\False}{\lif}[1]}, checked
      [\sFmla{\False}{(\formula{A} \lif \formula{B})},
        just={\TRule{\False}{\lif}[1]}
        [\sFmla{\True}{\lnot \formula{A}}, just={\TRule{\True}{\lor}[2]}]
        [\sFmla{\True}{\formula{B}}, just={\TRule{\True}{\lor}[2]}]
      ]
    ]
  ]
\end{oltableau}
Nous n'avons pas encore appliqué $\TRule{\False}{\lif}$ à la ligne~$3$ ;
faisons-le maintenant sur les deux branches. Rappelons que nous ne
cochons !!a{signed formula} qu'après avoir appliqué la règle
correspondante sur toutes les branches ouvertes. Il est donc utile
d'appliquer une règle à l'extrémité de chaque branche contenant la
!!{signed formula} visée. Nous n'aurons ainsi pas à revenir à cette
!!{signed formula} plus bas sur les différentes branches.
\begin{oltableau}
  [\sFmla{\False}{(\lnot \formula{A} \lor \formula{B}) \lif
      (\formula{A} \lif \formula{B})}, just=\TAss, checked
    [\sFmla{\True}{\lnot \formula{A} \lor \formula{B}},
      just={\TRule{\False}{\lif}[1]}, checked
      [\sFmla{\False}{(\formula{A} \lif \formula{B})},
        just={\TRule{\False}{\lif}[1]}, checked
        [\sFmla{\True}{\lnot \formula{A}}, just={\TRule{\True}{\lor}[2]}
          [\sFmla{\True}{\formula{A}}, just={\TRule{\False}{\lif}[3]}
            [\sFmla{\False}{\formula{B}}, just={\TRule{\False}{\lif}[3]}]
          ]
        ]
        [\sFmla{\True}{\formula{B}}, just={\TRule{\True}{\lor}[2]}
          [\sFmla{\True}{\formula{A}}, just={\TRule{\False}{\lif}[3]}
            [\sFmla{\False}{\formula{B}}, just={\TRule{\False}{\lif}[3]}, close]
          ]
        ]
      ]
    ]
  ]
\end{oltableau}
La branche de droite est maintenant fermée. Sur celle de gauche, nous
pouvons encore appliquer $\TRule{\True}{\lnot}$ à la ligne~$4$. On
obtient alors~$\sFmla{\False}{!A}$, ce qui ferme la branche gauche :
\begin{oltableau}
  [\sFmla{\False}{(\lnot \formula{A} \lor \formula{B}) \lif
      (\formula{A} \lif \formula{B})}, just=\TAss, checked
    [\sFmla{\True}{\lnot \formula{A} \lor \formula{B}},
      just={\TRule{\False}{\lif}[1]}, checked
      [\sFmla{\False}{(\formula{A} \lif \formula{B})},
        just={\TRule{\False}{\lif}[1]}, checked
        [\sFmla{\True}{\lnot \formula{A}}, just={\TRule{\True}{\lor}[2]}
          [\sFmla{\True}{\formula{A}}, just={\TRule{\False}{\lif}[3]}
            [\sFmla{\False}{\formula{B}}, just={\TRule{\False}{\lif}[3]}
              [\sFmla{\False}{\formula{A}},
                just={\TRule{\True}{\lnot}[4]}, close]
            ]
          ]
        ]
        [\sFmla{\True}{\formula{B}}, just={\TRule{\True}{\lor}[2]}
          [\sFmla{\True}{\formula{A}}, just={\TRule{\False}{\lif}[3]}
            [\sFmla{\False}{\formula{B}},
              just={\TRule{\False}{\lif}[3]}, close]
          ]
        ]
      ]
    ]
  ]
\end{oltableau}
\end{ex}

\begin{ex}
On peut construire des !!{tableau}s ayant un nombre quelconque de
!!{signed formula}s pour hypothèses. Il faut souvent appliquer plusieurs
règles qui produisent des ramifications, et !!a{tableau} peut en général
comporter un nombre quelconque de branches. Considérons par exemple
!!a{tableau} pour
$\{\sFmla{\True}{!A \lor (!B \land !C)}, \sFmla{\False}{(!A \lor !B) \land (!A
  \lor !C)}\}$. Commençons par appliquer $\TRule{\True}{\lor}$ à la
première hypothèse :
\begin{oltableau}
  [\sFmla{\True}{\formula{A} \lor (\formula{B} \land \formula{C})},
    just=\TAss, checked
    [\sFmla{\False}{(\formula{A} \lor \formula{B}) \land
        (\formula{A} \lor \formula{C})}, just=\TAss
      [\sFmla{\True}{\formula{A}}, just={\TRule{\True}{\lor}[1]}]
      [\sFmla{\True}{\formula{B} \land \formula{C}},
        just={\TRule{\True}{\lor}[1]}]
    ]
  ]
\end{oltableau}
Nous pouvons maintenant appliquer $\TRule{\False}{\land}$ à la
ligne~$2$. Faisons-le simultanément sur les deux branches ; nous pouvons
donc cocher la ligne~$2$ :
\begin{oltableau}
  [\sFmla{\True}{\formula{A} \lor (\formula{B} \land \formula{C})},
    just=\TAss, checked
    [\sFmla{\False}{(\formula{A} \lor \formula{B}) \land
        (\formula{A} \lor \formula{C})}, just=\TAss, checked
      [\sFmla{\True}{\formula{A}}, just={\TRule{\True}{\lor}[1]}
        [\sFmla{\False}{\formula{A} \lor \formula{B}},
          just={\TRule{\False}{\land}[2]}]
        [\sFmla{\False}{\formula{A} \lor \formula{C}},
          just={\TRule{\False}{\land}[2]}]
      ]
      [\sFmla{\True}{\formula{B} \land \formula{C}},
        just={\TRule{\True}{\lor}[1]}
        [\sFmla{\False}{\formula{A} \lor \formula{B}},
          just={\TRule{\False}{\land}[2]}]
        [\sFmla{\False}{\formula{A} \lor \formula{C}},
          just={\TRule{\False}{\land}[2]}]
      ]
    ]
  ]
\end{oltableau}
Appliquons ensuite $\TRule{\False}{\lor}$ sur toutes les branches qui
contiennent $!A \lor !B$ :
\begin{oltableau}
  [\sFmla{\True}{\formula{A} \lor (\formula{B} \land \formula{C})},
    just=\TAss, checked
    [\sFmla{\False}{(\formula{A} \lor \formula{B}) \land
        (\formula{A} \lor \formula{C})}, just=\TAss, checked
      [\sFmla{\True}{\formula{A}}, just={\TRule{\True}{\lor}[1]}
        [\sFmla{\False}{\formula{A} \lor \formula{B}},
          just={\TRule{\False}{\land}[2]}, checked
          [\sFmla{\False}{\formula{A}}, just={\TRule{\False}{\lor}[4]}
            [\sFmla{\False}{\formula{B}},
              just={\TRule{\False}{\lor}[4]}, close]
          ]
        ]
        [\sFmla{\False}{\formula{A} \lor \formula{C}},
          just={\TRule{\False}{\land}[2]}]
      ]
      [\sFmla{\True}{\formula{B} \land \formula{C}},
        just={\TRule{\True}{\lor}[1]}
        [\sFmla{\False}{\formula{A} \lor \formula{B}},
          just={\TRule{\False}{\land}[2]}, checked
          [\sFmla{\False}{\formula{A}}, just={\TRule{\False}{\lor}[4]}
            [\sFmla{\False}{\formula{B}},
              just={\TRule{\False}{\lor}[4]}]
          ]
        ]
        [\sFmla{\False}{\formula{A} \lor \formula{C}},
          just={\TRule{\False}{\land}[2]}]
      ]
    ]
  ]
\end{oltableau}
La branche située tout à gauche est désormais fermée. Appliquons
maintenant $\TRule{\False}{\lor}$ à $!A \lor !C$ :
\begin{oltableau}
  [\sFmla{\True}{\formula{A} \lor (\formula{B} \land \formula{C})},
    just=\TAss, checked
    [\sFmla{\False}{(\formula{A} \lor \formula{B}) \land
        (\formula{A} \lor \formula{C})}, just=\TAss, checked
      [\sFmla{\True}{\formula{A}}, just={\TRule{\True}{\lor}[1]}
        [\sFmla{\False}{\formula{A} \lor \formula{B}},
          just={\TRule{\False}{\land}[2]}, checked
          [\sFmla{\False}{\formula{A}},
            just={\TRule{\False}{\lor}[4]}
            [\sFmla{\False}{\formula{B}},
              just={\TRule{\False}{\lor}[4]}, close]
          ]
        ]
        [\sFmla{\False}{\formula{A} \lor \formula{C}},
          just={\TRule{\False}{\land}[2]}, checked
          [\sFmla{\False}{\formula{A}},
            just={\TRule{\False}{\lor}[4]}, move by=2
            [\sFmla{\False}{\formula{C}},
              just={\TRule{\False}{\lor}[4]}, close]
          ]
        ]
      ]
      [\sFmla{\True}{\formula{B} \land \formula{C}},
        just={\TRule{\True}{\lor}[1]}
        [\sFmla{\False}{\formula{A} \lor \formula{B}},
          just={\TRule{\False}{\land}[2]}, checked
          [\sFmla{\False}{\formula{A}},
            just={\TRule{\False}{\lor}[4]}
            [\sFmla{\False}{\formula{B}},
              just={\TRule{\False}{\lor}[4]}
            ]
          ]
        ]
        [\sFmla{\False}{\formula{A} \lor \formula{C}},
          just={\TRule{\False}{\land}[2]}, checked
          [\sFmla{\False}{\formula{A}},
            just={\TRule{\False}{\lor}[4]}, move by=2
            [\sFmla{\False}{\formula{C}},
              just={\TRule{\False}{\lor}[4]}]
          ]
        ]
      ]
    ]
  ]
\end{oltableau}
Pour plus de clarté, nous avons placé plus bas le résultat de la seconde
application de $\TRule{\False}{\lor}$. Ce déplacement n'était pas
nécessaire ici, puisque les justifications auraient été identiques.

Deux branches restent ouvertes, et $\sFmla{\True}{!B \land !C}$ à la
ligne~$3$ n'est toujours pas cochée. Appliquons-lui
$\TRule{\True}{\land}$ pour obtenir un !!{tableau} fermé :
\begin{oltableau}
  [\sFmla{\True}{\formula{A} \lor (\formula{B} \land \formula{C})},
    just=\TAss, checked
    [\sFmla{\False}{(\formula{A} \lor \formula{B}) \land
        (\formula{A} \lor \formula{C})}, just=\TAss, checked
      [\sFmla{\True}{\formula{A}}, just={\TRule{\True}{\lor}[1]}
        [\sFmla{\False}{\formula{A} \lor \formula{B}},
          just={\TRule{\False}{\land}[2]}, checked
          [\sFmla{\False}{\formula{A}},
            just={\TRule{\False}{\lor}[4]}
            [\sFmla{\False}{\formula{B}},
              just={\TRule{\False}{\lor}[4]}, close]
          ]
        ]
        [\sFmla{\False}{\formula{A} \lor \formula{C}},
          just={\TRule{\False}{\land}[2]}, checked
          [\sFmla{\False}{\formula{A}},
            just={\TRule{\False}{\lor}[4]}
            [\sFmla{\False}{\formula{C}},
              just={\TRule{\False}{\lor}[4]}, close]
          ]
        ]
      ]
      [\sFmla{\True}{\formula{B} \land \formula{C}},
        just={\TRule{\True}{\lor}[1]}, checked
        [\sFmla{\False}{\formula{A} \lor \formula{B}},
          just={\TRule{\False}{\land}[2]}, checked
          [\sFmla{\False}{\formula{A}},
            just={\TRule{\False}{\lor}[4]}
            [\sFmla{\False}{\formula{B}},
              just={\TRule{\False}{\lor}[4]}
              [\sFmla{\True}{\formula{B}},
                just={\TRule{\True}{\land}[3]}
                [\sFmla{\True}{\formula{C}},
                  just={\TRule{\True}{\land}[3]},close
                ]
              ]
            ]
          ]
        ]
        [\sFmla{\False}{\formula{A} \lor \formula{C}},
          just={\TRule{\False}{\land}[2]}, checked
          [\sFmla{\False}{\formula{A}},
            just={\TRule{\False}{\lor}[4]}
            [\sFmla{\False}{\formula{C}},
              just={\TRule{\False}{\lor}[4]}
              [\sFmla{\True}{\formula{B}},
                just={\TRule{\True}{\land}[3]}
                [\sFmla{\True}{\formula{C}},
                  just={\TRule{\True}{\land}[3]},close
                ]
              ]
            ]
          ]
        ]
      ]
    ]
  ]
\end{oltableau}

À titre de comparaison, voici un !!{tableau} fermé pour le même ensemble
d'hypothèses, dans lequel les règles sont appliquées dans un autre ordre :
\begin{oltableau}
  [\sFmla{\True}{\formula{A} \lor (\formula{B} \land \formula{C})},
    just=\TAss, checked
    [\sFmla{\False}{(\formula{A} \lor \formula{B}) \land
        (\formula{A} \lor \formula{C})}, just=\TAss, checked
      [\sFmla{\False}{\formula{A} \lor \formula{B}},
        just={\TRule{\False}{\land}[2]}, checked
        [\sFmla{\False}{\formula{A}},
          just={\TRule{\False}{\lor}[3]}
          [\sFmla{\False}{\formula{B}},
            just={\TRule{\False}{\lor}[3]}
            [\sFmla{\True}{\formula{A}},
              just={\TRule{\True}{\lor}[1]},close
            ]
            [\sFmla{\True}{\formula{B} \land \formula{C}},
              just={\TRule{\True}{\lor}[1]}, checked
              [\sFmla{\True}{\formula{B}},
                just={\TRule{\True}{\land}[6]}
                [\sFmla{\True}{\formula{C}},
                  just={\TRule{\True}{\land}[6]},close
                ]
              ]
            ]
          ]
        ]
      ]
      [\sFmla{\False}{\formula{A} \lor \formula{C}},
        just={\TRule{\False}{\land}[2]}, checked
        [\sFmla{\False}{\formula{A}},
          just={\TRule{\False}{\lor}[3]}
          [\sFmla{\False}{\formula{C}},
            just={\TRule{\False}{\lor}[3]}
            [\sFmla{\True}{\formula{A}},
              just={\TRule{\True}{\lor}[1]},close
            ]
            [\sFmla{\True}{\formula{B} \land \formula{C}},
              just={\TRule{\True}{\lor}[1]}, checked
              [\sFmla{\True}{\formula{B}},
                just={\TRule{\True}{\land}[6]}
                [\sFmla{\True}{\formula{C}},
                  just={\TRule{\True}{\land}[6]},close
                ]
              ]
            ]
          ]
        ]
      ]
    ]
  ]
\end{oltableau}
\end{ex}

\begin{prob}
Construire un !!{tableau} fermé pour chacun des ensembles suivants :
\begin{enumerate}
\item $\sFmla{\True}{!A \land (!B \land !C)}, \sFmla{\False}{(!A \land !B) \land !C}$.
\item $\sFmla{\True}{!A \lor (!B \lor !C)}, \sFmla{\False}{(!A \lor !B) \lor !C}$.
\item $\sFmla{\True}{!A \lif (!B \lif !C)}, \sFmla{\False}{!B \lif (!A \lif !C)}$.
\item $\sFmla{\True}{!A}, \sFmla{\False}{\lnot\lnot !A}$.
\end{enumerate}
\end{prob}

\begin{prob}
Construire un !!{tableau} fermé pour chacun des ensembles suivants :
\begin{enumerate}
\item $\sFmla{\True}{(!A \lor !B) \lif !C}, \sFmla{\False}{!A \lif !C}$.
\item $\sFmla{\True}{(!A \lif !C) \land (!B \lif !C)}, \sFmla{\False}{(!A \lor !B) \lif !C}$.
\item $\sFmla{\False}{\lnot(!A \land \lnot !A)}$.
\item $\sFmla{\True}{!B \lif !A}, \sFmla{\False}{\lnot !A \lif \lnot !B}$.
\item $\sFmla{\False}{(!A \lif \lnot !A) \lif \lnot !A}$.
\item $\sFmla{\False}{\lnot(!A \lif !B) \lif \lnot !B}$.
\item $\sFmla{\True}{!A \lif !C}, \sFmla{\False}{\lnot (!A \land \lnot !C)}$.
\item $\sFmla{\True}{!A \land \lnot !C}, \sFmla{\False}{\lnot (!A \lif !C)}$.
\item $\sFmla{\True}{!A \lor !B}, \sFmla{\True}{\lnot !B},
  \sFmla{\False}{!A}$.
\item $\sFmla{\True}{\lnot !A \lor \lnot !B}, \sFmla{\False}{\lnot(!A \land !B)}$.
\item $\sFmla{\False}{(\lnot !A \land \lnot !B) \lif\lnot(!A \lor !B)}$.
\item $\sFmla{\False}{\lnot(!A \lor !B) \lif (\lnot !A \land \lnot !B)}$.
\end{enumerate}
\end{prob}

\begin{editorial}
Note de l'édition française : dans le neuvième item du groupe précédent,
l'original regroupe par erreur $!A \lor !B$ et $\lnot !B$ sous un seul
signe~$\True$. Ils sont rétablis comme deux formules signées distinctes,
conformément aux exercices parallèles de déduction naturelle et de
calcul des séquents.
\end{editorial}

\begin{prob}
Construire un !!{tableau} fermé pour chacun des ensembles suivants :
\begin{enumerate}
\item $\sFmla{\True}{\lnot(!A \lif !B)}, \sFmla{\False}{!A}$.
\item $\sFmla{\True}{\lnot(!A \land !B)}, \sFmla{\False}{\lnot !A \lor \lnot !B}$.
\item $\sFmla{\True}{!A \lif !B}, \sFmla{\False}{\lnot !A \lor !B}$.
\item $\sFmla{\False}{\lnot \lnot !A \lif !A}$.
\item $\sFmla{\True}{!A \lif !B}, \sFmla{\True}{\lnot !A \lif !B}, \sFmla{\False}{!B}$.
\item $\sFmla{\True}{(!A \land !B) \lif !C}, \sFmla{\False}{(!A \lif !C) \lor (!B \lif !C)}$.
\item $\sFmla{\True}{(!A \lif !B) \lif !A}, \sFmla{\False}{!A}$.
\item $\sFmla{\False}{(!A \lif !B) \lor (!B \lif !C)}$.
\end{enumerate}
\end{prob}

\end{document}
